\subsection{Fields of representatives}

DEFINITION 1. --- Let A be a local ring. A field of representatives of A is any subfield K of A such that the canonical homomorphism of A onto \(\kappa_{\mathrm{A}}\) induces an isomorphism of K onto \(\kappa_{\mathrm{A}}\) (in other words, such that \(\mathrm{A} = \mathrm{K} + \mathfrak{m}_{\mathrm{A}}\) ).

A field of representatives of A can exist only if A has a characteristic. This condition is sufficient when A is separated and complete. More precisely, we have the following theorem:

THEOREM 1. --- Let A be a separated and complete local ring of characteristic p.

\begin{enumerate}
\def\labelenumi{\alph{enumi})}
\item
  Suppose \(p = 0\) and let \((x_{\lambda})_{\lambda \in \Lambda}\) be a family of elements of A whose classes modulo \(\mathfrak{m}_{\mathbf{A}}\) form a transcendence basis of \(\kappa_{\mathbf{A}}\) over \(\mathbf{Q}\) . There exists a unique field of representatives of A containing the elements \(x_{\lambda}\) .
\item
  Suppose \(p \neq 0\) . Let \((x_{\lambda})_{\lambda \in \Lambda}\) be a family of elements of A whose classes modulo \(\mathfrak{m}_{\mathbf{A}}\) form a p-basis of \(\kappa_{\mathbf{A}}\) (A, V, p. 95). There exists a unique field of representatives of A containing the elements \(x_{\lambda}\) . It is a Cohen subring of A.
\end{enumerate}

Suppose we have \(p = 0\) so that A is a Q-algebra. Every transcendence basis of \(\kappa_{\mathrm{A}}\) over Q being separating, the assertion \(a)\) follows from prop. 1 of no. 1 applied to the case \(k_0 = \mathbf{Q}\) , \(\mathbf{K} = \kappa_{\mathrm{A}}\) .

Suppose now that we have \(p \neq 0\) . Then one has \(p1_{\mathbf{A}} = 0\) , and every Cohen subring C of A satisfies \(p\mathbf{C} = 0\) . In other words, the notions of field of representatives and Cohen subring of A coincide. The assertion \(b)\) then follows from §2, no. 2, th. 1.

COROLLARY 1. --- Let A be a separated and complete local ring, whose residue field is an algebraic extension of Q. There then exists a unique field of representatives of A. Indeed, the ring A is of characteristic 0 (no. 1).

COROLLARY 2. --- Let A be a local ring, separated and complete, of characteristic \(p \neq 0\) . Suppose that the residue field \(\kappa_{\mathrm{A}}\) is perfect. Then there exists a unique field of representatives of A, namely the set of multiplicative representatives.

Cor. 2 follows immediately from th. 1 and prop. 7 of § 2, no. 4.

THEOREM 2. --- Let A be a complete Noetherian local ring of dimension d containing a field. Let K be a field of representatives of A, and let m be the dimension of the vector space \(\mathfrak{m}_{\mathrm{A}} / \mathfrak{m}_{\mathrm{A}}^{2}\) over the field K.

\begin{enumerate}
\def\labelenumi{\alph{enumi})}
\item
  There exists an ideal \(\mathfrak{a}\) of \(\mathbf{K}[[\mathrm{T}_1,\dots,\mathrm{T}_m]]\) such that the K-algebra A is isomorphic to \(\mathbf{K}[[\mathrm{T}_1,\dots,\mathrm{T}_m]] / \mathfrak{a}\) .
\item
  There exists a sub-K-algebra A' of A, isomorphic to \(K[[T_{1}, \ldots, T_{d}]]\) and such that A is a finite algebra over A'.
\item
  Suppose the Noetherian local ring A is regular, i.e. \(d = m\) . Then there exists a K-isomorphism from A onto \(\mathbf{K}[[\mathbf{T}_1, ..., \mathbf{T}_d]]\) .
\end{enumerate}

Let \(t_{1},\ldots,t_{m}\) be elements of \(\mathfrak{m}_{\mathrm{A}}\) whose classes modulo \(\mathfrak{m}_{\mathrm{A}}^{2}\) generate the K-vector space \(\mathfrak{m}_{\mathrm{A}}/\mathfrak{m}_{\mathrm{A}}^{2}\) . By lemma 3 of § 2, no. 5, there exists a surjective K-homomorphism from \(\mathrm{K}[[\mathrm{T}_{1},\ldots,\mathrm{T}_{m}]]\) to A, mapping \(\mathrm{T}_{i}\) to \(t_{i}\) for \(1\leqslant i\leqslant m\) . This proves \(a)\) .

Similarly, the assertion \(b\) ) follows from Lemma 4 of loc. cit. and from the existence of a maximal secant sequence for A (VIII, § 3, no. 2, th. 1).

Finally, assertion c) is none other than Cor. 3 of Th. 1 of VIII, § 5, no. 2.

